\begin{abstract}
Although large-scale unsupervised language models (LMs) are capable of acquiring extensive world knowledge and developing some reasoning abilities, exerting fine-grained control over their outputs remains challenging because of the inherently unsupervised nature of their training regime. Current approaches to enhance the controllability of these models typically involve collecting human evaluations that compare the relative quality of model outputs, and subsequently adapting the unsupervised LM to match these preferences—commonly through reinforcement learning from human feedback (RLHF). However, RLHF introduces significant complexity and instability; it involves first learning a reward model that embodies human preferences, and then using reinforcement learning to tune the LM so as to maximize this learned reward, all the while maintaining proximity to the initial model parameters. In this work, we propose a novel formulation for the reward model within RLHF that enables the direct retrieval of the optimal policy via a closed-form solution, thus transforming the standard RLHF problem into one solvable with a straightforward classification loss. Our proposed method, named \textit{Direct Preference Optimization} (DPO), is robust, efficient, and demands significantly less computational effort, as it removes the necessity for LM sampling during fine-tuning and does not require extensive hyperparameter search. Through empirical evaluation, we demonstrate that DPO can adapt LMs to human preferences on par with, or surpassing, prior approaches. In particular, DPO outperforms RLHF methods based on PPO in sentiment control, and achieves equal or superior results in tasks such as summarization and single-turn dialogue, while greatly simplifying the training and implementation process.
\end{abstract}

\section{Introduction}
Large-scale unsupervised language models (LMs) trained on vast corpora have demonstrated remarkable and often unexpected capabilities~\citep{chowdhery2022palm, brown2020language, touvron2023llama,bubeck2023sparks}. Yet, these LMs are built upon datasets created by humans possessing divergent goals, varying priorities, and heterogeneous levels of expertise. Not all human objectives or proficiencies encoded in such data are suitable or beneficial for the systems to replicate. For instance, in the case of an AI coding assistant, while it is important for the model to \textit{recognize} prevalent programming errors in order to address them effectively, we would nonetheless prefer that it \emph{emulate} the instances of superior coding ability—often a minority presence in its data. In the same vein, a model's \textit{awareness} of a widespread misconception (e.g., one believed by half of people) is useful diagnostically, but we do not want the system to propagate this misconception in half of its outputs. Thus, it is vital to \emph{select} from among a model’s broad-ranging \textit{knowledge and competencies} the particular responses and behavioral patterns that are aligned with our intended use—an imperative for achieving safe, effective, and controllable AI~\citep{ouyang2022training}. Whereas common techniques employ reinforcement learning (RL) to guide LMs according to human preferences, we demonstrate that the RL-driven objective currently used in practice can actually be solved exactly with a straightforward binary cross-entropy formulation, leading to a significant simplification of the overall preference alignment workflow.

Typically, to impart preferred behavioral tendencies, present-day methods utilize human-constructed preference datasets that exemplify the sorts of conduct deemed safe and advantageous. This preference alignment occurs after a foundational phase of unsupervised pre-training conducted over extensive textual data. While one naive strategy for preference shaping involves supervised fine-tuning on exemplary human-authored responses, state-of-the-art procedures center on reinforcement learning from human (or model-generated) feedback—referred to as RLHF or RLAIF~\citep{christiano2017deep,bai2022constitutional}. In these frameworks, a reward model is trained to replicate human judgments, which then serves as the basis for RL optimization of the policy underlying the LM. The policy is tuned to generate outputs that maximize the modeled reward, all while preventing substantial divergence from its pre-trained version. Despite the superior performance achieved through RLHF, particularly in domains such as dialogue and code generation, these methods are substantially more involved than conventional supervised approaches. They require multiple rounds of LM training, active policy sampling during learning, and thus incur substantial computational expenditure.

In this work, we demonstrate a method for aligning language models with human preferences without relying on explicit reward models or reinforcement learning. We introduce \textit{{Direct Preference Optimization} (DPO)}, an algorithm designed to implicitly target the same objective as popular RLHF techniques—namely, maximizing reward under a KL-divergence regularization—while maintaining ease of implementation and simplicity in training. Conceptually, DPO modifies the relative log probabilities in favor of preferred responses over disfavored ones, but crucially, it utilizes a per-example adaptive importance weight to circumvent the model degradation issues observed when optimizing a naive probability ratio. Similar to current approaches, DPO is underpinned by a theoretical preference framework, such as the Bradley-Terry model \cite{bradley1952rankanalysis}, which quantifies the agreement between a reward function and observed preference data. Distinctly, instead of leveraging this model to derive a loss for reward model training followed by policy optimization against the learned reward, DPO reformulates the preference loss to directly depend on the policy using a change of variables. As a result, with a set of human preference judgments over model outputs, DPO can train a policy through a straightforward binary cross entropy loss, directly learning the optimal policy for an implicit reward shaped by the preference annotations.

The principal contribution of our study is the introduction of Direct Preference Optimization (DPO), a reinforcement learning–free approach for preference-based language model training. Empirical results indicate that DPO matches or surpasses state-of-the-art methods, including PPO-based RLHF, for preference learning in tasks such as sentiment control, summarization, and conversational modeling, using models up to 6B parameters in scale.

\section{Related Work}

Self-supervised language models, particularly as their scale increases, are capable of solving certain tasks in a zero-shot setting \citep{radford2019language} or by leveraging few-shot prompting strategies \citep{gpt3,megatron,chowdhery2022palm}. Nonetheless, the accuracy on downstream tasks and alignment with user intentions can be greatly enhanced by further fine-tuning these models using datasets comprised of instructions paired with human-generated completions \citep{mishra-etal-2022-cross,sanh2022multitask,chung2022scaling,thoppilan2022lamda}. This approach, commonly referred to as ‘instruction-tuning,’ allows large language models to extend their generalization capabilities beyond the explicit instruction-tuning examples and overall improves their practical utility \citep{chung2022scaling}. While instruction tuning has been effective, collecting human judgments that compare responses is generally more tractable than curating expert-level demonstrations. Therefore, more recent work has focused on fine-tuning large language models using datasets derived from human preference rankings, which has led to gains in areas such as translation \citep{kreutzer-etal-2018-reliability}, summarization \citep{stiennon2022learning,ziegler2020finetuning}, story generation \citep{ziegler2020finetuning}, and following instructions \citep{ouyang2022training,ramamurthy2023is}. The prevailing paradigm in these works involves first training a neural network reward model to represent preferences according to a framework such as the Bradley-Terry model \citep{bradley1952rankanalysis}, followed by adjusting a language model’s parameters via reinforcement learning to maximize the modeled reward. Common algorithms for this step include REINFORCE \citep{williams1992reinforce}, proximal policy optimization (PPO; \cite{schulman2017proximal}), and their variations \citep{ramamurthy2023is}. Related research has utilized instruction-following LLMs, which have been fine-tuned using human feedback, as generators of additional synthetic preference data targeted towards specific qualities like safety or harmlessness \citep{bai2022constitutional}, employing minimal human intervention through rubric-based textual guidelines for LLM output annotation. This emerging field synthesizes advancements in reinforcement learning-based language model optimization across diverse objectives~\citep{Ranzato2015SequenceLT,paulus2018a,wu2018learning} with more general approaches for extracting learning signals from human preference data \citep{christiano2017deep,kupcsik2018learning}. However, the process of fine-tuning large language models via reinforcement learning using relative preference data presents significant technical and practical difficulties; the present work introduces a theoretically-grounded method for preference optimization that obviates the need for reinforcement learning.

Beyond language-specific scenarios, the problem of learning policies based on preferences has been investigated within both the bandit and reinforcement learning paradigms, yielding several proposed methodologies. In particular, contextual bandit algorithms that utilize preferences or action rankings instead of explicit reward signals are referred to as contextual dueling bandits (CDB; \cite{yue2012karmed,dudik2015contextual}). In this framework, where absolute rewards are unavailable, theoretical analyses of CDBs replace the classic optimal policy concept with the \textit{von Neumann winner}: this is a policy whose average likelihood of prevailing over any alternative policy meets or exceeds 50\% \citep{dudik2015contextual}. Notably, CDBs typically receive preference feedback in an online fashion; by contrast, preference learning from humans often proceeds using a static dataset of previously annotated preference pairs \citep{yan2022human}. In a related vein, \textit{preference-based RL} (PbRL) employs binary preferences emerging from an \textit{unknown} evaluation or ‘scoring’ function rather than from reward signals \citep{BusaFekete2014,ruiz2023dueling}. There are various PbRL algorithms—including those that support leveraging offline preference data—but most approaches consist of two sequential steps: first, inferring the hidden scoring or reward function, and then optimizing policies against it \citep{jain2013learning,BusaFekete2014,christiano2017deep,sadigh2017active,kupcsik2018learning}. In this work, we advocate for a unified, end-to-end policy learning procedure that aims to optimize a policy in a manner directly consistent with preference data, bypassing explicit reward model estimation.

\section{Preliminaries}\label{section:prelims}

We summarize the RLHF procedure described by \citeauthor{ziegler2020finetuning} and further developed in works such as \citep{stiennon2022learning, bai2022training, ouyang2022training}. The standard process is comprised of three distinct stages: (1) supervised fine-tuning (SFT), (2) acquisition of preference data and training of a reward model, and (3) reinforcement learning for policy optimization.

\textbf{SFT}: The pipeline commences with supervised fine-tuning of a pre-trained language model, utilizing curated, high-quality datasets tailored to the intended application (for example, dialogue or summarization). This results in the fine-tuned model $\pi^\text{SFT}$.

\textbf{Reward Modelling Phase}: Subsequently, the SFT model is conditioned on prompts $x$ and produces pairs of candidate responses $(y_1, y_2)\sim \pi^\text{SFT}(y \mid x)$. Human annotators then express a preference between these alternatives, formally denoted $y_w\succ y_l \mid x$, where $y_w$ and $y_l$ correspond to the favored and disfavored responses, respectively, among $(y_1, y_2)$. These preference labels are understood to be the output of some unobserved, latent reward function $r^*(y, x)$. Multiple statistical models exist to represent such preference structures, with the Bradley-Terry (BT) \cite{bradley1952rankanalysis} model being a commonly used approach. More general models such as Plackett-Luce \citep{plackett1975analysis, luce2012individual} are applicable as well in cases where a ranked list of answers is available. The BT framework characterizes the human preference probability distribution $p^*$ as follows:

\begin{equation}\label{eq:bradley-terry}
    p^*(y_1\succ y_2 \mid x)=\frac{\exp\left(r^*(x, y_1)\right)}{\exp\left(r^*(x, y_1)\right) + \exp\left(r^*(x, y_2)\right)}.
\end{equation}

Given a fixed dataset of comparison tuples, $\mathcal{D} = \bigl\{x^{(i)}, y_w^{(i)}, y_l^{(i)}\bigr\}_{i=1}^N$, where samples are drawn according to $p^*$, one can introduce a parameterized reward function $r_{\phi}(x, y)$ and fit its parameters by maximizing likelihood. Interpreting the problem as a binary classification task, the negative log-likelihood can be formalized as:

\begin{equation}\label{eq:reward_model}
    \mathcal{L}_R(r_{\phi}, \mathcal{D}) = -\mathbb{E}_{(x, y_w, y_l)\sim \mathcal{D}}\bigl[\log \sigma(r_{\phi}(x, y_w)- r_{\phi}(x, y_l))\bigr]
\end{equation}

where $\sigma$ denotes the logistic sigmoid. In language model applications, $r_{\phi}(x, y)$ is typically initialized with parameters from the SFT policy $\pi^\text{SFT}(y \mid x)$, augmented by a linear head attached to the last transformer hidden state, thereby outputting a scalar reward estimate~\cite{ziegler2020finetuning}. To ensure stability and minimize variance in the reward values, prior work commonly normalizes the rewards so that $\mathbb{E}_{x,y\sim \mathcal{D}}\left[r_\phi(x, y)\right] = 0$ across $x$.

\textbf{RL Fine-Tuning Phase}: In the reinforcement learning phase, feedback to the language model is delivered through the learned reward signal. Adhering to previous methodologies~\citep{jaques2017sequence, jaques2020human}, the training objective can be written as

\begin{equation}\label{eq:RL}
\max_{\pi_{\theta}}  \mathbb{E}_{x\sim \mathcal{D}, y\sim \pi_{\theta}(y \mid x)}\bigl[r_{\phi}(x, y)\bigr] - \beta\mathbb{D}_{\textrm{KL}}\bigl[\pi_{\theta}(y\mid x)\mid \mid \pi_\text{ref}(y\mid x)\bigr],
\end{equation}

with $\beta$ moderating the degree of divergence allowed from the reference policy $\pi_\text{ref}$, i.e., the SFT policy $\pi^\text{SFT}$. Commonly, the RL policy $\pi_\theta$ is also initialized using $\pi^\text{SFT}$. Incorporating the KL regularization term is crucial; it constrains the model from departing excessively from the distribution for which the reward model provides reliable assessments, thus preserving output diversity and averting collapse to only a few reward-maximizing sequences. Owing to the discrete nature of text, gradients for this objective are unavailable; as a result, reinforcement learning algorithms are applied. Standard practice~\citep{ziegler2020finetuning, stiennon2022learning, bai2022training, ouyang2022training} involves defining the reward as ${r(x, y) = r_{\phi}(x, y) -\beta (\log \pi_{\theta}(y\mid x) - \log \pi_\text{ref}(y\mid x))}$, and employing PPO~\cite{schulman2017proximal} for optimization. 

\section{Direct Preference Optimization}\label{sec:DPO}

Recognizing the computational and algorithmic obstacles when scaling RL to large language models, we seek a more direct technique for policy improvement using preference data. Distinct from prior RLHF pipelines that first estimate a reward then optimize it, our strategy is built upon a reward model formulation that permits analytic derivation of the optimal policy, thereby eliminating the need for a reinforcement learning stage. As detailed below, the central insight is to utilize the direct analytical correspondence between reward functions and their optimal policies. This allows us to reframe loss functions that typically operate over rewards into ones defined over policies instead. Through this transformation, our method avoids learning an explicit and independent reward estimator, yet retains the benefits of modeling human preferences (e.g., via the Bradley-Terry framework). Thus, the

\textbf{Deriving the DPO objective.} We begin with the reinforcement learning objective previously introduced in Eq.~\ref{eq:RL}, which employs a general reward function $r$. In line with established work~\citep{peters2007reinforcement, peng2019advantage, korbak2022reinforcement, go2023aligning}, one can readily demonstrate that the policy that optimally maximizes the reward subject to a KL constraint, as defined in Eq.~\ref{eq:RL}, is given by:

\begin{equation}\label{eq:op_policy}
    \pi_r(y\mid x) = \frac{1}{Z(x)}\pi_\text{ref}(y\mid x)\exp\left(\frac{1}{\beta}r(x, y)\right),
\end{equation}

where the normalization constant $Z(x) = \sum_{y}\pi_\text{ref}(y\mid x)\exp\left(\frac{1}{\beta}r(x, y)\right)$ acts as a partition function. For the detailed mathematical derivation, see Appendix \ref{app:derivation1}. Despite the availability of an MLE-based approximation $r_{\phi}$ for the underlying reward function $r^*$, directly calculating the partition function $Z(x)$ remains computationally challenging~\citep{korbak2022reinforcement, go2023aligning}. This limitation complicates the application of the policy form in practice. Nonetheless, it is possible to rearrange Eq.~\ref{eq:op_policy} to represent the reward in terms of the optimal policy $\pi_r$, the reference policy $\pi_\text{ref}$, and $Z(\cdot)$. By applying a logarithm to both sides of Eq.~\ref{eq:op_policy} and conducting algebraic manipulations, we have:

\begin{equation}\label{eq:main_eq}
    r(x,y) =\beta \log \frac{\pi_r(y\mid x)}{\pi_\text{ref}(y\mid x)} + \beta \log Z(x).
\end{equation}

This formulation is equally applicable to the true reward $r^*$ and its associated optimal policy $\pi^*$. Notably, the Bradley-Terry model characterizes preferences using only reward differences between completions: ${p^*(y_1 \succ y_2 \mid x) = \sigma(r^*(x, y_1) - r^*(x, y_2))}$. Substituting the expression from Eq.~\ref{eq:main_eq} for $r^*(x,y)$ into the Bradley-Terry preference formulation (Eq.~\ref{eq:bradley-terry}), the terms involving the partition function cancel each other. Consequently, the likelihood of a human preference depends solely on the optimal policy $\pi^*$ and the reference policy $\pi_\text{ref}$. Therefore, under the Bradley-Terry framework, the optimal RLHF policy $\pi^*$ is characterized by the following preference model:

\begin{equation}\label{eq:objective}
    p^*(y_1\succ y_2 \mid x)=\frac{1}{1 + \exp\left(\beta \log \frac{\pi^*(y_2\mid x)}{\pi_\text{ref}(y_2\mid x)} - \beta \log \frac{\pi^*(y_1\mid x)}{\pi_\text{ref}(y_1\mid x)}\right)}
\end{equation}

A detailed derivation is provided in Appendix~\ref{app:derivation2}. Although Eq.~\ref{eq:objective} is formulated using the Bradley-Terry model, a similar approach extends to the more general Plackett-Luce models~\citep{plackett1975analysis, luce2012individual}, as discussed in Appendix~\ref{app:plackett_luce_models}.

With the probability of human preference now expressed as a function of the optimal policy rather than a reward function, it is possible to define a maximum likelihood objective for a parameterized policy $\pi_\theta$. Drawing an analogy to reward modeling (see Eq.~\ref{eq:reward_model}), the policy learning objective becomes:

\begin{equation}\label{eq:optimum_model}
    \mathcal{L}_\text{DPO}(\pi_{\theta}; \pi_\text{ref}) = -\mathbb{E}_{(x, y_w, y_l)\sim \mathcal{D}}\left[\log \sigma \left(\beta \log \frac{\pi_{\theta}(y_w\mid x)}{\pi_\text{ref

In this manner, we learn an implicit reward via a distinct parameterization, wherein the resulting optimal policy corresponds directly to $\pi_\theta$. Additionally, because our approach aligns with the reparametrized Bradley-Terry model, it inherits certain theoretical guarantees, such as consistency provided that the preference data satisfy suitable distributional assumptions \cite{bong2022generalized}. Section~\ref{sec:theory} delves further into the theoretical analysis of DPO and compares its properties to related methodologies.

\textbf{Mechanistic interpretation of the DPO update.} To gain deeper insight into the operation of DPO, it is instructive to inspect the gradient of its loss, $\mathcal{L}_\text{DPO}$. Specifically, the parameter gradient is given by:

\begin{multline*}\label{eq:gradient}
    \nabla_\theta \mathcal{L}_\text{DPO}(\pi_\theta;\pi_\text{ref}) = \\ -\beta\mathbb{E}_{(x, y_w, y_l) \sim \mathcal{D}} \bigg[\underbrace{\sigma(\hat{r}_\theta(x, y_l) - \hat{r}_\theta (x, y_w))}_\text{emphasizes higher loss where reward ranks are incorrect}\bigg[\underbrace{\nabla_\theta\log \pi(y_w \mid x)}_\text{pushes up the probability of $y_w$} - \underbrace{\nabla_\theta\log\pi(y_l \mid x)}_\text{pushes down the probability of $y_l$}\bigg]\bigg],
\end{multline*}

where the implicit reward, $\hat{r}_\theta(x, y) = \beta \log \frac{\pi_\theta(y \mid x)}{\pi_\text{ref}(y \mid x)}$, is defined through the language model $\pi_\theta$ and the reference $\pi_\text{ref}$ (discussed further in Section~\ref{sec:theory}). Conceptually, this loss gradient serves to favorably increase the probability of preferred responses $y_w$ and decrease the probability for less preferred ones $y_l$. Crucially, the contribution of each sample is scaled according to the degree to which $\hat{r}_\theta$ mistakenly favors the dispreferred completion, modulated by $\beta$, thus reflecting the severity of the error in the implicit reward's ordering as well as the impact of the KL constraint. Empirically, this scaling factor proves to be vital; omitting it (i.e., using an unweighted variant) can result in model collapse, as shown in Appendix Table~\ref{tab:unlikelihood_generations}.

\textbf{DPO overview.} The typical DPO workflow consists of the following steps: 1) For each input prompt $x$, generate two candidate completions $y_1, y_2 \sim \pi_\text{ref}(\cdot \mid x)$, annotate these with human preference labels to form the offline preference dataset $\mathcal{D} = \{x^{(i)}, y_w^{(i)}, y_l^{(i)}\}_{i=1}^N$, and 2) update the language model $\pi_\theta$ by minimizing $\mathcal{L}_\text{DPO}$, given $\pi_\text{ref}$, the dataset $\mathcal{D}$, and the desired $\beta$ value. 

In real-world applications, it is preferable to utilize publicly available preference datasets rather than creating new samples and soliciting fresh human feedback. When such datasets have been created using $\pi^\text{SFT}$, we set the reference policy as $\pi_\text{ref} = \pi^\text{SFT}$ if possible. On the other hand, in the absence of $\pi^\text{SFT}$, $\pi_\text{ref}$ is initialized by maximizing the likelihood over preferred completions ${(x, y_w)}$, i.e., ${\pi_\text{ref} = \argmax_{\pi}\mathbb{E}_{x, y_w \sim \mathcal{D}}\left[\log \pi(y_w \mid x)\right]}$. This initialization strategy helps reduce distributional discrepancies between the inaccessible true reference distribution and the $\pi_\text{ref}$ actually employed in DPO. Additional information on implementation choices and hyperparameter selections can be found in Appendix~\ref{app:implementation}.

\section{Theoretical Analysis of DPO}

Here, we provide a deeper examination of the DPO framework, supplying theoretical justifications and connecting its strengths to known shortcomings of actor-critic methods used in RLHF, such as PPO~\cite{schulman2017proximal}.

\label{sec:theory}

\subsection{Your Language Model Is Secretly a Reward Model} 

DPO enables learning policies without explicitly training a reward model or performing reinforcement learning, instead leveraging a single maximum likelihood training objective. It is important to observe that the DPO objective in Eq. \ref{eq:main_eq} matches a Bradley-Terry model where rewards are parametrized by $r^*(x, y) = \beta \log\frac{\pi^*_\theta(y \mid x)}{\pi_\text{ref}(y \mid x)}$. The parametric policy $\pi_{\theta}$ is trained analogously to optimizing the reward model as in Eq. \ref{eq:reward_model} under a variable transformation. In what follows, we develop the formal foundations for this reparameterization, demonstrate that it does not restrict the expressivity of possible learned reward functions, and establish that the optimal policy can be exactly reconstructed. We initiate this discussion by introducing an equivalence relationship among reward functions.

\begin{definition}
Two reward functions $r(x, y)$ and $r'(x, y)$ are defined to be equivalent if and only if there exists some function $f$ such that ${r(x, y)-r'(x, y) = f(x)}$.
\end{definition}

This notion of equivalence establishes a true equivalence relation among reward functions, resulting in a partitioning of the set of all such functions into distinct equivalence classes. This observation leads naturally to the following lemmas:

\begin{lemma}\label{lemma:same_prefrence} Within the Plackett-Luce framework—including the Bradley-Terry model—any pair of reward functions belonging to the same equivalence class give rise to identical preference distributions.
\end{lemma}

\begin{lemma}\label{lemma:same_policy}
    Any two reward functions from a common equivalence class yield the same optimal policy when considering the constrained RL formulation.
\end{lemma}

The demonstrations of these lemmas are direct, so we postpone their details to Appendix \ref{app:lemma1}. The first lemma formalizes a well-established aspect of under-specification found in the Plackett-Luce model family \cite{plackett1975analysis}. Because of this ambiguity, identifiability constraints are frequently introduced to guarantee any consistency for MLE estimators in Eq. \ref{eq:reward_model} \cite{bong2022generalized}. The second lemma clarifies that all members of a given equivalence class produce the same optimal policy; thus, our interest, for the purpose of final policy optimization, is limited to recovering some member from the correct equivalence class. The main result below is established in Appendix~\ref{app:thm1}:

\begin{theorem}\label{thm:main}
    Assuming mild regularity conditions, any reward class that matches the Plackett-Luce (or specifically Bradley-Terry) criteria admits a representation of the form ${r(x, y) = \beta \log \frac{\pi(y\mid x)}{\pi_\text{ref}(y\mid x)}}$, where $\pi(y\mid x)$ is an appropriate model and $\pi_\text{ref}(y \mid x)$ is a prescribed reference model.
\end{theorem}

\begin{sproof}
    Given a reward function $r(x, y)$, let $\pi_r(y \mid x)$ be the corresponding optimal model as defined by Eq. \ref{eq:op_policy}. We claim that every reward function within the equivalence class of $r$ is representable by the aforementioned reparameterization. To see this, consider the mapping
\begin{equation}
    f(r; \pi_\text{ref}, \beta)(x, y) = r(x, y) - \beta\log\sum_{y}\pi_\text{ref}(y\mid x)\exp\left(\frac{1}{\beta}r(x, y)\right)
\end{equation}
This operator normalizes the original reward by subtracting a log-partition function, which depends only on the input $x$; thus, $f(r; \pi_\text{ref}, \beta)(x, y)$ lies in the same equivalence class as $r(x, y)$. Substituting the expression for $r$ from Eq.~\ref{eq:main_eq}, which universally holds, yields $f(r; \pi_\text{ref}, \beta)(x, y) = \beta \log \frac{\pi_r(y\mid x)}{\pi_\text{ref}(y\mid x)}$. Therefore, the mapping $f$ constructs a representative from the equivalence class of $r$ with the desired structural form, ensuring that the proposed reparameterization does not restrict the generality of the reward function.
\end{sproof}

Theorem~\ref{thm:main} can also be interpreted as identifying, for each equivalence class, the particular reward function selected by the DPO reparameterization. Specifically, this reward function must satisfy:

\begin{equation}\label{eq:lag_p}
     \sum_{y}\underbrace{\pi_\text{ref}(y\mid x)\exp\left(\frac{1}{\beta}r(x, y)\right)}_{=\pi(y\mid x)\text{, using Thm.~\ref{thm:main} reparam.}} = 1,
\end{equation}

which ensures that $\pi(y\mid x)$ constitutes a proper probability distribution, meaning it is non-negative and sums to unity. Examining Eq.~\ref{eq:op_policy}, one sees that Eq.~\ref{eq:lag_p} serves as the partition function for the optimal policy arising from the reward $r(x, y)$. The central observation in the DPO algorithm is that by enforcing specific constraints on the under-determined Plackett-Luce (and particularly the Bradley-Terry) family of preference models, one can both retain the expressiveness of the reward model class and, at the same time, guarantee that the optimal policy in Eq.~\ref{eq:op_policy} remains analytically tractable for every input $x$.

\subsection{Instability of Actor-Critic Algorithms} Our framework further enables the analysis of instability issues encountered by standard actor-critic methods employed in RLHF workflows, such as PPO. Here, we examine the RL fine-tuning phase described in Section \ref{section:prelims}. Connections can be established with the control as inference perspective \cite{levine2018reinforcement}, especially in the context of the constrained RL problem stated in \ref{eq:RL}. Given a model family $\pi_{\theta}(y\mid x)$, we consider minimizing the divergence $\mathbb{D}_{\text{KL}}[\pi_{\theta}(y|x) \mid \mid \pi^*(y\mid x)]$, where $\pi^*$ refers to the optimal policy as given by Eq.~\ref{eq:optimum_model} for the reward $r_{\phi}(y, x)$. Manipulating this objective algebraically yields:

\begin{equation}\label{eq:AC}
    \max_{\pi_{\theta}}\mathbb{E}_{\pi_{\theta}(y\mid x)}\bigg[\underbrace{r_{\phi}(x, y) -\beta\log\sum_{y}\pi_\text{ref}(y\mid x)\exp\left(\frac{1}{\beta}r_{\phi}(x, y)\right)}_{f(r_{\phi}, \pi_\text{ref}, \beta)} - \underbrace{\beta\log\frac{\pi_{\theta}(y\mid x)}{\pi_\text{ref}(y\mid x)}}_{\text{KL}}\bigg]
\end{equation}

The objective here aligns with that optimized in earlier studies 
\citep{ziegler2020finetuning, stiennon2022learning, bai2022training, ouyang2022training}, employing the DPO-equivalent reward associated with the reward class $r_{\phi}$. In this framework, the normalization component within $f(r_{\phi}, \pi_\text{ref}, \beta)$ can be seen as representing the soft value function of the reference policy $\pi_\text{ref}$. Although this normalization does not impact the optimality of the solution, omitting it may lead to high variance in the policy gradient, thereby destabilizing training. One approach to include this normalization is to introduce a trainable value function, but this method can be challenging to optimize in practice. Another strategy, previously adopted, is to normalize rewards relative to a baseline provided by a human completion, effectively treating it as a Monte Carlo estimate based on a single sample for the normalization term. By contrast, the DPO reparameterization introduces a reward structure that obviates the need for such baselines.

\section{Experiments}
Here, we present an empirical assessment of DPO's capacity to train policies solely using preference feedback. We begin in a controlled text-generation context and investigate DPO’s sample efficiency with respect to balancing reward maximization and KL-divergence minimization against the reference policy, benchmarking it against commonly used preference learning methods, including PPO. Subsequently, we test DPO’s efficacy on tasks involving larger models and greater RLHF complexity, such as dialogue and summarization. Our results demonstrate that DPO typically matches or surpasses robust baselines like PPO-based RLHF and outperforms the strategy of choosing the best trajectory among $N$ samples using a learned reward, all while requiring minimal hyperparameter adjustment. Details about the experimental protocol are provided before we report these results; further information is available in Appendix~\ref{app:exp_details}.

\textbf{Tasks.} We conduct experiments across three distinct open-ended text generation tasks. In all cases, the algorithms are trained to learn a policy from a preference dataset $\mathcal{D}=\bigl\{x^{(i)}, y_w^{(i)}, y_l^{(i)}\bigr\}_{i=1}^N$. The first task, \textbf{controlled sentiment generation}, utilizes $x$ as a prefix from an IMDb movie review \cite{maas-EtAl:2011:ACL-HLT2011}. Here, the goal for the policy is to generate an output $y$ with a positive sentiment. To ensure the evaluation remains controlled, we construct preference pairs over generated completions by employing a pre-trained sentiment classifier, ensuring that $p(\text{positive}\mid x,y_w)>p(\text{positive}\mid x,y_l)$. For SFT, GPT-2-large is fine-tuned on the training portion of the IMDb reviews until performance plateaus (details provided in App~\ref{app:sentiment_details}). In the \textbf{summarization} task, $x$ represents a Reddit forum post, and the task is to produce a summary $y$ that captures the main points. Adopting the methodology of earlier studies, we leverage the Reddit TL;DR summarization dataset \citep{volske-etal-2017-tl} together with human preference data curated by \citeauthor{stiennon2022learning}. SFT for this task relies on a model fine-tuned on reference summaries authored by humans,\footnote{\url{https://huggingface.co/CarperAI/openai_summarize_tldr_sft}} and we use the TRLX framework \citep{leandro_von_werra_2023_7790115} to implement RLHF. The collection of preference annotations was performed by \citeauthor{stiennon2022learning}, sampled from outputs of a distinct, but similarly fine-tuned, SFT model. The final task, \textbf{single-turn dialogue}, considers $x$ as an arbitrary human query—ranging from scientific inquiries to personal advice. Here, the policy's aim is to generate an informative and engaging response $y$; the dataset used is Anthropic's Helpful and Harmless dialogue corpus \citep{bai2022training}, comprising 170k interactions between human users and a conversational assistant. Each dialogue ends with two model-generated responses and a human preference annotation. Because no specialized SFT model is available in this context, we derive the SFT model by fine-tuning a generic pre-trained language model solely on the responses preferred by annotators.

\textbf{Evaluation.} We utilize two primary evaluation strategies in our experimental setup. For the controlled sentiment generation scenario, each algorithm’s ability to optimize the constrained reward maximization objective is examined by plotting the achievable trade-off between reward and KL-divergence with respect to the reference policy. This is feasible because we possess the ground-truth reward function (a sentiment classifier). In contrast, when evaluating in real-world scenarios where the reward function is unknown, we adopt the metric of \textit{win rate} relative to a baseline policy. Here, GPT-4 acts as a stand-in for human annotators to judge summary quality and the helpfulness of responses in summarization and single-turn dialogue tasks, respectively. In summarization tasks, the baseline comprises reference summaries from the test set, while in dialogue, we compare against the responses labeled as preferred in the dataset. Given prior evidence indicating that LMs are sometimes more reliable automatic evaluators than established metrics \citep{Chen2023ExploringTU}, we further perform a human evaluation (see Sec.~\ref{sec:human-judgments}) to validate the suitability of GPT-4 as a proxy evaluator. Our findings show that GPT-4’s assessments are strongly aligned with human judgments, often yielding human-GPT-4 agreement rates as high as or higher than those among human annotators.

\textbf{Methods.} Beyond DPO, we benchmark several alternative techniques for aligning language models to human preferences. For summarization, we apply zero-shot prompting to \textbf{GPT-J} \citep{gpt-j}, while for dialogue, we employ 2-shot prompting with \textbf{Pythia-2.8B} \citep{biderman2023pythia}. We also include the \textbf{SFT} model and a variant, \textbf{Preferred-FT}, which is trained via supervised fine-tuning using the selected response $y_w$—sourced either from the SFT output (in sentiment control and summarization) or from a generic LM (in single-turn dialogue). The \textbf{Unlikelihood} method~\citep{welleck2019neural} constitutes another pseudo-supervised approach, encouraging the model to increase the likelihood of $y_w$ while concurrently decreasing that of $y_l$, controlled by an optional weighting factor $\alpha\in[0,1]$ for the penalization term. Additionally, we assess \textbf{PPO} \citep{schulman2017proximal}, which leverages a reward model inferred from human preference data, and \textbf{PPO-GT}, an idealized variant that operates with the actual reward function in the sentiment-controlled environment. Our sentiment experiments employ both an out-of-the-box PPO-GT implementation \cite{leandro_von_werra_2023_7790115} and an enhanced version with reward normalization and refined hyperparameters, the latter modifications also applied to standard PPO with learned rewards. As a strong baseline, we also evaluate the \textbf{Best of $N$} approach, which involves generating $N$ samples from the SFT model (or Preferred-FT for dialogue), and selecting the top-scoring sample per a reward model derived from preference annotations. Although highly effective, this method is computationally expensive, as it necessitates generating $N$ completions per query during evaluation.

\subsection{How well can DPO optimize the RLHF objective?}

The reward maximization objective in RLHF, subject to a KL constraint, aims to find a balance between maximizing reward and preventing the learned policy from straying significantly from the reference policy. Consequently, when evaluating and comparing different approaches, it is essential to jointly consider the reward attained and the degree of KL divergence; higher rewards are less meaningful if achieved at the cost of substantially increased KL. Figure~\ref{fig:frontier-tldr-main} displays the trade-off between reward and KL divergence (the reward-KL frontier) for several algorithms in the sentiment domain. For each method, we carry out several independent training runs, each with a distinct setting for policy conservativeness (specifically, PPO's target KL $\in\{3,6,9,12\}$; $\beta \in \{0.05,0.1,1,5\}$ and $\alpha\in\{0.05,0.1,0.5,1\}$ for unlikelihood; and different random seeds for preferred-FT), resulting in a total of 22 training runs. At intervals of 100 training steps, and continuing until the algorithm converges, we assess the current policy's performance by evaluating on a held-out set of test prompts. We measure the mean reward according to the true reward function, as well as the average sequence-level KL,\footnote{Namely, the sum of KL-divergences at each time step in the sequence.} computed against the reference policy $\text{KL}\left(\pi\mid \mid \pi_\text{ref}\right)$. Our findings show that DPO achieves by far the most favorable reward-KL trade-off, realizing the greatest reward for comparatively small KL divergence. This outcome is particularly significant for two main reasons. First, despite both DPO and PPO being designed to optimize the same underlying objective, DPO consistently attains a more efficient trade-off, with DPO's performance on the reward/KL frontier unequivocally surpassing that of PPO. Second, even when PPO is granted access to ground truth reward information (PPO-GT), DPO outperforms PPO and defines a superior reward-KL frontier.

\subsection{Can DPO scale to real preference datasets?}
\label{sec:dpo-real-datasets}
We proceed by assessing the effectiveness of DPO fine-tuning in the context of summarization and single-turn dialogue tasks. In the domain of summarization, widely used automatic metrics like ROUGE often show a weak correlation with human judgment~\citep{stiennon2022learning}, and previous studies have shown that using PPO for fine-tuning language models with human preference data can yield more desirable summaries. To benchmark various approaches, we generate samples on the TL;DR summarization dataset’s test split and calculate the mean win rate compared to test set references. For each method, completions are drawn using temperature values ranging between 0.0 and 1.0, and the corresponding win rates are illustrated in Figure~\ref{fig:frontier-tldr-main} (right). All models, including DPO, PPO, and Preferred-FT, are fine-tuned from the same GPT-J SFT model\footnote{\url{https://huggingface.co/CarperAI/openai_summarize_tldr_sft}}. Results indicate that DPO achieves a win rate of roughly 61\% at temperature 0.0, which outperforms PPO’s best win rate of about 57\% at the same temperature. Additionally, DPO surpasses the $N$-best baseline in terms of its peak win rate. It's important to highlight that DPO’s $\beta$ parameter was not thoroughly optimized in our experiments, suggesting that actual DPO performance could be higher. DPO also demonstrates notably greater resilience to variations in sampling temperature, whereas PPO’s effectiveness drops sharply at higher temperatures, reaching the level of the unmodified GPT-J model. Preferred-FT, on the other hand, shows limited improvement over the SFT baseline. Furthermore, we conduct a direct human evaluation between DPO and PPO, described in Section~\ref{sec:human-judgments}, where DPO samples drawn at temperature 0.25 are chosen 58\% of the time over PPO samples generated at the same temperature.

For single-turn dialogue experiments, we assess the various approaches using the portion of the test split from the Anthropic HH dataset \citep{bai2022training} that contains a single exchange between a human and assistant. In the GPT-4 evaluations, the preferred completions from the test set are treated as references to calculate the win rate across different methods. Since there is no widely-accepted SFT model for this setting, our procedure starts with the pre-trained Pythia-2.8B model. We then apply Preferred-FT to train a reference model on selected completions, ensuring the outputs stay within the model’s typical distribution, followed by DPO training. For benchmarking, we include a comparison against the best result among 128 Preferred-FT completions (we observed that the performance of the Best of $N$ baseline stabilizes at 128 samples for this dataset; refer to Appendix Figure~\ref{fig:best-of-n}) and a version of the Pythia-2.8B base model with 2-shot prompting, noting that DPO matches or surpasses the strongest temperatures of other methods. Additionally, we analyze an RLHF model trained with PPO on the Anthropic HH dataset \footnote{\url{https://huggingface.co/reciprocate/ppo_hh_pythia-6B}} obtained from a reputable repository \footnote{\url{https://github.com/CarperAI/trlx/tree/main/examples/hh}}, but could not identify a combination of prompt and temperature that consistently outperformed the vanilla Pythia-2.8B. Drawing from our TL;DR results and acknowledging that both Best of 128 and PPO target the same reward signal, we treat Best of 128 as an approximate stand-in for PPO-level outcomes. To summarize, DPO stands out as the only method with practical computational efficiency that improves upon the Anthropic HH dataset’s preferred completions, and it achieves comparable or superior results to the costly Best of 128 baseline. Lastly, Figure~\ref{fig:dialogue-main} illustrates that DPO attains its optimal performance in relatively few training iterations.

\subsection{Generalization to a new input distribution}

To further assess how PPO and DPO perform under distributional changes, we tested the policies obtained from our Reddit TL;DR summarization task on a distinct domain: news articles found in the test set of the CNN/DailyMail dataset \citep{nallapati-etal-2016-abstractive}, utilizing the optimal sampling temperatures previously identified for TL;DR (0 and 0.25). The performance metrics are shown in Table~\ref{tab:ood}. We measured the GPT-4 win rate relative to the ground-truth news article summaries, employing the same GPT-4 (C) prompt as in the Reddit TL;DR task but swapping ``forum post'' with ``news article.'' On this new input domain, DPO substantially exceeds the PPO policy in performance. These results offer preliminary evidence that DPO policies maintain strong generalization comparable to PPO, even though DPO does not leverage the extra unlabeled Reddit TL;DR prompts used in PPO training.

\subsection{Validating GPT-4 judgments with human judgments}
\label{sec:human-judgments}
To validate the consistency of GPT-4's scoring with human preferences, we ran a user study using outputs from the TL;DR summarization experiment and two alternate GPT-4 prompts. The \textbf{GPT-4 (S)} (simple) prompt requests an evaluation of which summary more effectively covers the crucial content from the post. The \textbf{GPT-4 (C)} (concise) prompt in addition asks which summary is more succinct; we selected this prompt after observing that, under the \textbf{GPT-4 (S)} prompt, GPT-4 shows a preference for verbose, repetitive summaries, unlike human raters. Complete versions of the prompts are in Appendix~\ref{app:prompts}. We conducted three comparisons—using the top-performing (DPO, temp. 0.25), the lowest-performing (PPO, temp. 1.0), and an intermediate (SFT, temp. 0.25) method, aiming to capture a range of summary qualities—each compared to PPO with greedy decoding (its most effective temperature). Our analysis shows that, for both prompts, GPT-4's preferences align with human judgments about as frequently as humans concur with one another, indicating that GPT-4 serves as a reliable stand-in for human evaluation (owing to limited human rater availability, we gather multiple human scores only for DPO and PPO-1 matchups). Notably, the \textbf{GPT-4 (C)} prompt produces win rates that are better aligned with human assessments; therefore, we adopt this prompt for our principal results in Section~\ref{sec:dpo-real-datasets}. Further methodological specifics about the human study, including details about the rating interface and participating raters, can be found in Appendix~\ref{app:human-study}.

\section{Discussion}
Preference-based learning offers an effective and scalable approach for developing capable and aligned language models. In this work, we proposed DPO, a straightforward training methodology that enables language models to learn from preferences without utilizing reinforcement learning techniques. Instead of reshaping the preference learning problem to fit conventional RL frameworks for compatibility with standard RL tools, DPO establishes a direct correspondence between language model policies and reward structures. This formulation allows for training language models to adhere to human preferences via a basic cross-entropy objective, eliminating the need for reinforcement learning while preserving generality. Remarkably, DPO achieves comparable or superior results to prevalent RLHF approaches, such as those employing PPO, and requires minimal hyperparameter tuning. As a result, DPO substantially simplifies the process of aligning language models with human preferences.

\textbf{Limitations \& Future Work.} Our findings bring to light a number of open questions that warrant further exploration. For instance, how robustly does a policy learned with DPO generalize to distributions beyond those encountered during training, particularly in comparison with policies trained using explicit reward signals? Preliminary evidence points to DPO's generalization abilities being on par with PPO-based methods, though extensive investigations remain necessary. Moreover, it is yet to be determined whether techniques such as self-labeling with the DPO policy can be leveraged to capitalize on unlabeled prompts effectively. Another avenue for research involves understanding the characteristics of reward over-optimization in the direct preference optimization context; for instance, does the minor performance drop observed in Figure~\ref{fig:dialogue-main}-right reflect such over-optimization? Furthermore, while our experiments consider models up to 6B parameters, scaling DPO to modern, much larger architectures is an intriguing area for future studies. In the context of evaluations, we observe that GPT-4’s calculated win rates are sensitive to prompt variations; optimizing automated evaluation methodologies could be a fruitful direction for future work. Finally, DPO holds promise for numerous applications outside human preference alignment in language models, including the training of generative models across other data modalities.

\end{document}